% ============================================================
%  三阶段实验计划 — "夹心饼干 → 苦瓜"结构改造（v2，2026-08-25）
%  本文件为纯正文，被 \input 进 menu.tex（主报告第 4 章 三阶段实验日志）
%  改动记录：v1 基础上并入精读改进——
%    ① 补 E 组（糖×预氧化交叉）与无糖对照（A0/B0，解耦 PVP 碳背景）
%    ② 新增 Zr 掺量轴（约5/20 mol%），对接 anand2020/liu2016 的保碳-耐温平衡
%    ③ 阻抗匹配判定先行（|Zin/Z0| 与 α），对接 liu2025/cheng2025
%    ④ 1250 °C 支线（激活碳热还原产气造孔）
%    ⑤ 预氧化时间梯度（0.5/1/2 h）、PAN 替换 PVP 第二轮对照（老师推荐）
%    ⑥ 表征：Raman 三参数（ωG+ID/IG+D′）、XRD 少研磨防 t→m、形貌验证先行、热导口径标注
%  v2.1（2026-09-02，与实验日志 CSV 对账后定稿）：
%    ① 硅烷比按历史实际执行定 5:1（MTMS:DMDMS = 25:5，非 1:1），75 g 罐 = 62.5+12.5 g
%    ② 外层 PVP 定 1.72 g/10 g 基液（20260615/0712 口径）
%    ③ Zr 0.64 g/10 g 基液 $\approx$ Zr:Si 5 mol%（anand2020 甜点）；历史 15+15+0.96 为 0.64 g/10 g 硅烷 $\approx$2.5 mol%，记录注明口径
%    ④ C-Zr↓ 定为 0.64 g（=标准，平行重复）；C-Zr↑ = 2.56 g $\approx$ 20 mol%
%    ⑤ 新增 D3 支线（B 配方、硅烷比 3:1），探低比可纺性与碳增益
% ============================================================

\section{背景与目标}

仿生参考：Yang 等 2026（\textit{Advanced Science}，DOI 10.1002/advs.76343）制备了苦瓜仿生结构的 C-SiC 纤维气凝胶——纤维\textbf{表面布满颗粒突起与沟槽}、\textbf{截面呈"致密外壳 + 疏松内核"}，一举同时实现超低热导（9.12 mW·m$^{-1}$·K$^{-1}$）与电磁屏蔽。其核壳是在 1600–1800 °C 下由 \ce{SiC_{x}O_{y}N_{z}} 分解释放 \ce{SiO}/\ce{CO} 气体、气体被表面碳层拦下发生碳热还原而\textbf{自发}形成的。

本计划的定位：\textbf{不复制其高温机理}（我们温度上限以 900 °C 为主，体系中 \ce{Zr} 会生成 \ce{ZrSiO4} 锁住 \ce{SiO2}、抑制碳热还原），而是用\textbf{工艺手段在较低温度内造出同样的形貌}："致密皮 + 疏松瓤 + 表面突起"。目标功能：吸波（表面多次散射、界面极化）与隔热（拉长热路、减重）双收。

老师推荐文献（Yang 2026）标注的两个亮点——\textbf{PAN 替换 PVP}、\textbf{离心纺丝}——作为本计划的后续探索方向（第二轮 P 系列，见下），本轮先做 PVP 体系的苦瓜结构验证。

\section{设计原理（三要点）}

\begin{enumerate}
  \item \textbf{表皮先固化}：空气预氧化（200–250 °C）使纤维表层 PVP 氧化交联、硅烷进一步缩合，形成"硬壳"；
  \item \textbf{内芯气化造孔}：惰性热解时内层 PVP 与葡萄糖分解产气（\ce{CO}/\ce{CO2}/焦油类），在内部留下疏松孔道；
  \item \textbf{气体外逸顶出突起}：被硬壳约束的气体从表层薄弱处逸出，形成颗粒突起与沟槽；升温速率控制突起大小。
\end{enumerate}

\section{配方}

\textbf{基液（基础液，先纯水预水解，可批量配好密封存放）}：三甲氧基硅烷(MTMS)与二甲氧基硅烷(DMDMS)按 \textbf{5:1 质量比}（25\,g:5\,g，与 2025-09 以来全部历史实验一致；批量 75\,g = 62.5\,g + 12.5\,g；D3 支线例外按 3:1，15\,g = 11.25\,g + 3.75\,g）混合成硅烷原液，取 10\,g，\textbf{加去离子水几滴（约 0.1--0.2\,mL，不加酸）}，室温搅拌\textbf{数小时至过夜}预水解（水被硅烷吸收成 Si--OH，形成交联位点），再与 DMF 10\,g 混匀（1:1）共 20\,g。硅烷混合液与基液均可提前批量配好、密封防潮存放（DMF 吸湿，用前检查无浑浊）。

\textbf{外层液（致密壳，外层不动、按现有工艺）}：取基液 10\,g，加 \textbf{PVP 1.72\,g}（现有外层配方定稿量，20260615/0712 口径），室温搅拌 6\,h；一个 mat 的两层外层按基液 20\,g + PVP 3.44\,g 配。

\textbf{外层液（PAN 版，第二轮探索）}：基液 10\,g + PAN（溶于 DMAc，用量初定 2\,g，需先做溶解性测试）——对应老师推荐文献的"PAN 替 PVP"。

\textbf{中层液（疏松瓤）}：取基液 10\,g，加 PVP 1.72\,g、正丙醇锆 0.64\,g（\textbf{0.64\,g/10\,g 基液 = Zr:Si $\approx$ 5 mol\%}，恰为 anand2020 保碳甜点；"C-Zr$\downarrow$" 即 0.64\,g 平行重复，C-Zr$\uparrow$ 用 2.56\,g $\approx$ 20 mol\%。\emph{口径提醒}：历史 15+15+0.96 基液是 0.64\,g/10\,g 硅烷 $\approx$ 2.5 mol\%，本计划在只含 5\,g 硅烷的 10\,g 基液上投 0.64\,g，实验本注明）、葡萄糖（0/2/4\,g）。葡萄糖先以去离子水预溶成糖水再加入（2\,g+3\,mL、4\,g+4--5\,mL 水，40--50 °C 加热至澄清，现配现用），室温搅拌 6\,h。

\textbf{加料顺序}：正丙醇锆忌游离水——先加 PVP 与基液搅匀，\textbf{再加正丙醇锆}（优先与预水解产生的 Si--OH 结合生成 Si--O--Zr），\textbf{最后加糖水}，加完尽快纺丝，缩短锆与游离水的接触。预水解采用纯水不加酸，避免残留酸在后续催化锆/葡萄糖水解。

\textbf{记录要求}：硅烷比 5:1（MTMS:DMDMS，D3 支线 3:1 例外）、基液（硅烷:DMF）1:1、Zr 0.64\,g/10\,g 基液 $\approx$ 5 mol\% 的口径在实验本写死；记录\textbf{实际投入的葡萄糖、去离子水与正丙醇锆质量}，不只记"几克"。

\section{样品组}

\textbf{主线（900 °C，8 组）}：

\begin{table}[htbp]
\centering
\caption{主线样品组：糖量 × 预氧化 × Zr 掺量}
\begin{tabularx}{\textwidth}{cccccX}
\toprule
样品 & 中层糖 & 空气预氧化 & Zr & 热解 & 目的 \\
\midrule
A   & 2\,g & 无       & 标准(0.64\,g) & 900 °C & 现有工艺基线 \\
A0  & 0\,g & 无       & 标准          & 900 °C & PVP 碳背景（无糖对照） \\
B   & 2\,g & 有（1\,h） & 标准         & 900 °C & 预氧化成壳 \\
B0  & 0\,g & 有（1\,h） & 标准         & 900 °C & 预氧化 × 无糖背景 \\
C   & 4\,g & 有（1\,h） & 标准         & 900 °C & 预氧化 + 高糖造孔 \\
E   & 4\,g & 无       & 标准          & 900 °C & 糖 × 预氧化交叉 \\
C-Zr↓ & 4\,g & 有（1\,h） & 0.64\,g（$\approx$5 mol\%） & 900 °C & 标准 Zr 重复验证 \\
C-Zr↑ & 4\,g & 有（1\,h） & 2.56\,g（$\approx$20 mol\%） & 900 °C & Zr 保碳上限（抑碳更强） \\
\bottomrule
\end{tabularx}
\end{table}

\textbf{支线}：

\begin{table}[htbp]
\centering
\caption{支线：温度、预氧化时间与 PAN 对照}
\begin{tabularx}{\textwidth}{lX}
\toprule
样品 & 说明 \\
\midrule
C-1250 / E-1250 & C、E 各取一份改烧 1250 °C（激活碳热还原产气造孔，衔接苦瓜机理；需原 1250 °C 炉子） \\
B-t0.5 / B-t2 & B 组预氧化时间梯度 0.5\,h / 2\,h（看壳厚与突起形貌，防中间层糖被烧掉） \\
D3 & B 配方、硅烷改 3:1，其余同 B \\
P 系列（第二轮） & 外层液 PVP$\rightarrow$PAN（溶 DMAc），先做溶解性与可纺性测试 \\
\bottomrule
\end{tabularx}
\end{table}

\textbf{D3 说明}：B 组配方（糖 2\,g、预氧化 1\,h、Zr 0.64\,g），仅硅烷原液比改 3:1——小批 15\,g = MTMS 11.25\,g + DMDMS 3.75\,g（加几滴水预水解过夜即可）。目的：低比可纺性与碳增益摸底，为下一轮整体切换探路；本轮主线一律 5:1，不混比。

\section{工艺步骤}

\begin{enumerate}
  \item 按样品表配液（外层液、中层液，含各对照组），搅拌 6\,h，记录实际葡萄糖/Zr 质量；
  \item 三层静电纺丝：外–中–外各 10\,mL，挤出速率 1.5\,mL/h，\textbf{连续纺织不中断}（硅油纸收集、全金属针头）；
  \item 揭下纺织物，记录初重；
  \item \textbf{空气预氧化}（仅 B/B0/C/C-Zr↓/C-Zr↑ 及支线 B-t0.5/B-t2/D3）：马弗炉空气，200–250 °C 保温按组（0.5/1/2\,h），使表皮固化；A/A0/E 跳过；
  \item 降温取出，转移至惰性气氛炉（先降温再切换气氛，两段必须分开）；
  \item \textbf{\ce{Ar} 热解}：200 °C 保温 1\,h 除残余溶剂 $\rightarrow$ 200–600 °C 慢速升温（气化段 1–2 °C/min）$\rightarrow$ 600 °C 至目标温度（主线 900 °C；C-1250/E-1250 为 1250 °C）保温 1\,h $\rightarrow$ 自然炉冷；
  \item \textbf{介电-匹配判定（先行）}：每组成型后先测复介电常数（2–18 GHz），按传输线理论计算 $|Z_{\mathrm{in}}/Z_0|$ 与衰减常数 $\alpha$——\textbf{阻抗匹配优先于衰减}（liu2025/cheng2025 结论）：若 $|Z_{\mathrm{in}}/Z_0|$ 远离 1，先调碳/糖或 Zr 掺量，再谈 RL 谷值；
  \item 记录终重与失重率，编号存档。
\end{enumerate}

\section{表征与判定}

\begin{table}[htbp]
\centering
\caption{表征手段与判定标准（形貌验证先行）}
\begin{tabularx}{\textwidth}{lX}
\toprule
表征 & 判定标准 \\
\midrule
SEM 表面（先做） & 预氧化组出现颗粒突起/沟槽；A 组表面光滑 \\
SEM 截面（先做） & 壳致密 + 芯疏松（核壳对比度） \\
BET 孔容/比表面 & C 组显著高于 A 组（疏松芯证据） \\
Raman（$\omega_{\mathrm{G}}+I_{\mathrm{D}}/I_{\mathrm{G}}+D^{\prime}$） & 自由碳状态与石墨化程度，\textbf{必须配同温无糖对照}（liu2016） \\
XRD & \textbf{少研磨/整束测}：研磨会诱发 t$\rightarrow$m \ce{ZrO2} 相变（9\%$\rightarrow$32\%） \\
热导率 & 标注口径（常温毡/高温/复合材料），勿混比 XieSong2019 常温毡值 \\
RL / EAB & 复介电算 $|Z_{\mathrm{in}}/Z_0|$、$\alpha$、RL、EAB；厚度 1–5 mm 全扫 \\
\bottomrule
\end{tabularx}
\end{table}

\section{注意事项}

\begin{itemize}
  \item \textbf{无糖对照必须做}（A0/B0）：PVP 本身是碳源，不控则 ID/IG、$\varepsilon''$ 变化无法归因；
  \item \textbf{阻抗匹配优先}：碳/导电相"适量最优"，不是越多越好（过量碳会阻抗失配、表面反射，见 cheng2025 的 ACET 25\% 反例）；
  \item \textbf{Zr 掺量是保碳-耐温的旋钮}：anand2020 显示 5 mol\% 保留碳最多、20 mol\% 抑碳最强——与"要碳吸波"目标取平衡；
  \item 预氧化温度 \textbf{不超过 250 °C}，预氧化时间别过长，防止把中间层葡萄糖也氧化烧掉（先 B-t0.5 验证）；
  \item 空气预氧化与 \ce{Ar} 热解两段气氛\textbf{必须分开}，先降温再切换；
  \item 葡萄糖吸湿、易析晶，\textbf{当天配液当天纺}，干燥后尽快烧结；
  \item 资源有限时\textbf{串行执行}：先 B（只验证预氧化成壳）$\rightarrow$ 再 C/E（糖与交叉）$\rightarrow$ 最后 Zr 系列；
  \item 第二轮 P 系列（PAN 替换 PVP）前先做溶解性/可纺性小试，PAN 需 DMAc 溶解并注意预氧化制度差异。
\end{itemize}
